% ID: SEL-D0081-OB-E02
% ESTADO: cotejada. Segunda pasada visual y borrador compilado revisados.
% Fuente DOC: Selectividad/sel_2012_matematicas_aplicadas/3 MATEMÁT. APLIC. CC. SS/EXAMEN 2-ANDALUCÍA.doc
% SHA-256 original: bcec57ae525fa5e23ce74631274a2abda0b5b16ac47ce6108ba8b5572d4aa634
% Página convertida: 2; opción B; ejercicio 2.

Sea $P(t)$ el porcentaje de células de un determinado tejido, afectadas por un cierto tipo de enfermedad transcurrido un tiempo $t$, medido en meses:
\[
P(t)=\begin{cases}
t^2&\text{si }0\leq t\leq5\\[0.5em]
\dfrac{100t-250}{t+5}&\text{si }t>5
\end{cases}.
\]
\begin{enumerate}[label=\alph*)]
\item \textbf{(0.5 puntos)} Estudie la continuidad de la función $P$.
\item \textbf{(0.75 puntos)} Estudie la derivabilidad de $P$ en $t=5$.
\item \textbf{(0.75 puntos)} Estudie la monotonía de dicha función e interprete la evolución del porcentaje de células afectadas.
\item \textbf{(0.5 puntos)} \textquestiondown{}En algún momento el porcentaje de células afectadas podría valer 50?
\end{enumerate}
